\section{Method}
\label{sec:method}

\noindent\textbf{Assumptions.}
Following the local-consistency assumption of
Voxgraph~\cite{reijgwart2020voxgraph}, we assume that TSDF fusion under the
drifted pre-correction trajectory remains locally valid: surface geometry and
connectivity are coherent in local neighborhoods although the map is globally
inconsistent. WarpVox corrects this global inconsistency after loop closure;
severely noisy or topologically corrupted source maps are outside its scope.

\subsection{Problem and Correction Inputs} Let $T^-(t),T^+(t)\in SE(3)$ map the sensor frame at time $t$ to the world frame before and after trajectory correction. The map built with $T^-$ consists of a TSDF $\Phi^-$, its weight field $w^-$, voxel size $s$, and truncation distance $\tau$. During online fusion, Voxblox continuously maintains the Marching-Cubes surface $\mathcal M^-=(V^-,F)$ of the active TSDF~\cite{10.1145/37402.37422}. On a trajectory correction, this mesh becomes the source surface. Let $\mathcal A$ denote the active TSDF cells that contribute triangles to $\mathcal M^-$. WarpVox stores the timestamped, modality-native range images of SLAM keyframes,
\begin{equation}
\mathcal C_{\rm KF}=\{k_i=(t_i,D_i):i\in{\rm KF}\},
\end{equation}
where ${\rm KF}$ indexes the retained SLAM keyframes and $D_i$ is an RGB-D depth image or a spherical LiDAR range image. Together with $T^-$, these measurements form the observation input to Stage~1. The measurements remain available across correction events. Fig.~\ref{fig:pipeline} summarizes the four stages.

\begin{figure}[t]
  \centering
  \includegraphics[width=\linewidth]{Figures/pipeline_v5.jpg}
  \caption{WarpVox extracts the current source surface from the fused TSDF, re-establishes projective correspondences with keyframe measurements, propagates the resulting targets through one proxy deformation per correction, and reconstructs a source-supported TSDF.}
  \label{fig:pipeline}
\end{figure}

\subsection{Stage 1: Keyframe--Surface Association}

Stage~1 projects the current surface into the keyframe range images and returns vertex-indexed observation records,
\begin{equation}
\begin{aligned}
\mathcal O&=\mathcal R(\mathcal M^-,T^-,\mathcal C_{\rm KF}),\\
\mathcal O_j&=\{(r,\alpha_{rj})\}.
\end{aligned}
\label{eq:resolved_records}
\end{equation}
where each record $r$ identifies a keyframe range image and a valid projected range sample, $t_r$ is its acquisition time, and $\alpha_{rj}\in[0,1]$ is the
effective reliability of associating it with source vertex $v_j^-$. For target construction, each vertex retains at most 15 association records, giving fixed-size per-vertex state. RGB-D restricts candidates with a camera-centre KD-tree, LiDAR with a
sensor-position grid and an 80\,m radius check. The correction
backend therefore receives vertex-indexed records and never re-fuses
measurements into the volume. Its contract is
\begin{equation}
(\Phi^-,w^-,\mathcal M^-,T^+,\mathcal O)
\longmapsto (\Phi^+,w^+),
\label{eq:contract}
\end{equation}
where the warped mesh is intermediate and $(\Phi^+,w^+)$ is a
native Voxblox TSDF volume.

For an RGB-D record $r$, pixel $\bar u_r=[u_r,v_r,1]^\top$ with measured
z-depth $z_r$ defines the invariant sensor-frame sample
\begin{equation}
m_r=z_rK^{-1}\bar u_r,
\qquad
\widetilde m_r=[m_r^\top,1]^\top.
\label{eq:sensor_sample}
\end{equation}
The inverse pre-correction pose projects source vertex $v_j^-$ into the record frame. Candidate agreement combines depth residual, three-dimensional position, and viewing direction. Let $n_j^-$ be the source normal, $\ell_{rj}^-$ the unit viewing direction in the pre-correction world frame, and $z_{rj}^{\rm proj}$ the projected vertex depth. We define
\begin{equation}
\begin{aligned}
e^g_{rj}
&=\left\|[T^-(t_r)\widetilde m_r]_{1:3}-v_j^-\right\|,\\
e^z_{rj}
&=z_{rj}^{\rm proj}-z_r,\\
g_{rj}
&=\left|(n_j^-)^\top \ell_{rj}^-\right|.
\end{aligned}
\label{eq:association_residuals}
\end{equation}
The effective association reliability is
\begin{equation}
\alpha_{rj}=
\exp\!\left(-\frac{e^g_{rj}}{0.048\,\mathrm{m}}\right)
\exp\!\left(-\frac{|e^z_{rj}|}{0.08\,\mathrm{m}}\right)
\frac{g_{rj}}{1+z_{rj}^{\rm proj}/8\,\mathrm{m}}.
\label{eq:association_weight}
\end{equation}
RGB-D candidates satisfy the depth-dependent geometric gate $e^g_{rj}\le 0.08\,\mathrm{m}+0.01\,z_{rj}^{\rm proj}$, and candidates with $\alpha_{rj}<0.05$ are rejected. LiDAR uses spherical projection with radial range $\rho$: candidates satisfy $e^g_{rj}\le 0.25\,\mathrm{m}+0.003\,\rho_{rj}^{\rm proj}$ and are scored by \begin{equation} \alpha_{rj}=\exp\!\left(-\frac{e^g_{rj}}{0.15\,\mathrm{m}}\right) \frac{g_{rj}}{1+\rho_{rj}^{\rm proj}/40\,\mathrm{m}}, \label{eq:association_weight_lidar} \end{equation} with no further $\alpha$ threshold. This establishes geometric support, not exact voxel-integration provenance. The attenuation in \eqref{eq:association_weight} ranks candidates during association; the separate factor in \eqref{eq:target_initial} models target reliability after association.

\subsection{Stage 2: Corrected Surface Targets}

After loop closure, the timestamp $t_r$ selects the corresponding corrected pose. Missing interior poses use linear translation interpolation and SLERP rotation over gaps of at most three times the median trajectory spacing~\cite{shoemake1985slerp}. The invariant sample then defines the corrected candidate
\begin{equation}
q_r^+=\left[T^+(t_r)\widetilde m_r\right]_{1:3}.
\label{eq:corrected_target}
\end{equation}
This candidate is geometrically consistent when the static-scene association, calibration, and
corrected pose are valid. Let
$\mathcal J_j=\{r:(r,\alpha_{rj})\in\mathcal O_j\}$. We first compute
\begin{equation}
\omega_{rj}=\frac{\alpha_{rj}}{1+z_r/3\,\mathrm{m}},
\qquad
\bar q_j=
\frac{\sum_{r\in\mathcal J_j}\omega_{rj}q_r^+}
     {\sum_{r\in\mathcal J_j}\omega_{rj}}.
\label{eq:target_initial}
\end{equation}
For LiDAR, $\omega_{rj}=\alpha_{rj}/(1+\|q_r^+-v_j^-\|/20\,\mathrm{m})$ and the floor in \eqref{eq:target_inliers} is $0.05\,\mathrm{m}$.
The reliability-weighted center is paired with an unweighted median residual
scale so that no single high-weight candidate controls both the estimate and
its rejection threshold:
\begin{equation}
\begin{aligned}
\theta_j&=\max\!\left(
0.02\,\mathrm{m},
2.5\operatorname*{median}_{r\in\mathcal J_j}
\|q_r^+-\bar q_j\|\right),\\
I_j&=\left\{r\in\mathcal J_j:
\|q_r^+-\bar q_j\|\le\theta_j\right\}.
\end{aligned}
\label{eq:target_inliers}
\end{equation}
When fewer than three candidates are available, $I_j=\mathcal J_j$. The final target is
\begin{equation}
v_j^*=
\frac{\sum_{r\in I_j}\omega_{rj}q_r^+}
     {\sum_{r\in I_j}\omega_{rj}}.
\label{eq:target}
\end{equation}
Each vertex with at least one valid record yields a target; the resulting controls are screened by the robust spatial-consistency filter of Stage~3.

\subsection{Stage 3: Surface Deformation}

The deformation stage propagates observation-derived motion; it does not
infer correction from shape alone. Let $d_j=v_j^*-v_j^-$ be a control
displacement. Before deformation, control-aware mesh cleaning preserves
connected components carrying observation-derived controls. A fixed
component-aware robust filter then removes isolated controls whose
displacements disagree with the surrounding accepted motion. Small
disconnected components retained only because they initially carried controls
are excluded if all such controls are subsequently rejected. Remaining
components with no surviving control inherit a robust local translation from
nearby accepted controls through synthetic constraints; these constraints
define otherwise underconstrained motion but are not treated as observations.

Component-separated spatial clustering~\cite{rossignac1993vertexclustering}
constructs a triangular proxy $\mathcal P^-=(P^-,F_P)$ with a target budget
of approximately 60K vertices and selects one QEM
representative~\cite{garland1997qem} per occupied cell. Let
$\pi:V^-\rightarrow P^-$ map source vertices to representatives and
$\mathcal K_a=\{j:\pi(j)=a\}$ collect their controls. For
$\mathcal K_a\neq\emptyset$, the initial proxy target is
\begin{equation}
p_a^{(0)}=p_a^-+
\frac{1}{|\mathcal K_a|}
\sum_{j\in\mathcal K_a}d_j.
\label{eq:proxy_target}
\end{equation}
A robust local-rigid consensus
~\cite{schoenemann1966procrustes,yao2020welsch} regularizes these proxy
targets. An observation-derived target is replaced by the fitted motion only
when a strict majority of the evaluated targets in the same connected
component are inconsistent with that fit; isolated disagreements remain
unchanged. Synthetic targets are regularized by the same local motion field.
Let $\mathcal C_P\subseteq\{1,\ldots,|P^-|\}$ index the resulting measured
or synthetic proxy targets $p_a^*$.

WarpVox performs one proxy ARAP solve per correction~\cite{ARAP_modeling:2007, 10.1007/978-3-662-05105-4_2,libigl}:
\begin{equation}
\begin{aligned}
\min_{P^+,\{R_a\in SO(3)\}}\quad
&\frac{1}{2}\sum_{a=1}^{|P^-|}
\sum_{b\in\mathcal N_P(a)}c_{ab}\\[-1mm]
&\left\|(p_a^+-p_b^+)-R_a(p_a^--p_b^-)\right\|^2,\\
\text{s.t.}\quad
&p_a^+=p_a^*,\qquad a\in\mathcal C_P .
\end{aligned}
\label{eq:arap}
\end{equation}
Here $\mathcal N_P(a)$ is the one-ring neighbourhood induced by $F_P$, and $c_{ab}=c_{ba}$ are symmetric cotangent-derived weights; the factor $1/2$ removes double counting. Once accepted, controls are imposed exactly so that ARAP does not attenuate the supplied trajectory correction. The standard local--global iterations solve this single optimization.

Proxy displacement is transferred to $V^-$ by topology-local barycentric
interpolation. Let $f_j$ be the closest triangle to $v_j^-$ in
$\operatorname{star}(\pi(j))$. With barycentric coordinates $\lambda_{jb}$ of
its closest point,
\begin{equation}
v_j^+=v_j^-+\sum_{b\in f_j}\lambda_{jb}(p_b^+-p_b^-),
\qquad \lambda_{jb}\ge0,\quad\sum_{b\in f_j}\lambda_{jb}=1.
\label{eq:proxy_transfer}
\end{equation}
The ARAP solve and topology-local transfer modify vertex positions only;
before the safety pass, they preserve the source connectivity and the
separation of disconnected components. A subsequent safety pass harmonically
relaxes regions exhibiting severe deformation-induced distortion. Faces whose
edge or area expansion remains unresolved are removed, but no new connections
are created and the surface is not retriangulated. The retained source and
warped meshes therefore share the same indexed face subset,
\begin{equation}
\bar{\mathcal M}^-=(\bar V^-,\bar F),\qquad
\mathcal M^+=(\bar V^+,\bar F),\qquad
\bar F\subseteq F.
\label{eq:retained_topology}
\end{equation}

\subsection{Stage 4: Source-Supported TSDF Reconstruction}

The warped surface supplies the corrected zero set
~\cite{sussman1999redistancing}, while the source TSDF supplies observed
support, sign orientation, integration weights, and free space. Let $n_b$ be
the number of voxels per block edge, $h=n_bs$ the block width, $o_b=hb$ the
origin of block $b$, and
\begin{equation}
\operatorname{vox}(b)=
\left\{o_b+s\left(k+\tfrac12\mathbf 1\right):
k\in\{0,\ldots,n_b-1\}^3\right\}
\label{eq:block_voxels}
\end{equation}
its voxel centres. Let $\mathcal B_{\rm fs}^+$ denote the additional blocks
reached by forward transport of observed source free space. With
$\mathcal N_1$ denoting a one-block Chebyshev dilation, the final allocation
domain is
\begin{equation}
\begin{aligned}
\mathcal B_V^+&=\{\lfloor \bar v/h\rfloor:\bar v\in\bar V^+\},\\
\mathcal B^+&=\mathcal N_1(\mathcal B_V^+)\cup\mathcal B_{\rm fs}^+,\\
\mathcal X^+&=\bigcup_{b\in\mathcal B^+}\operatorname{vox}(b).
\end{aligned}
\label{eq:allocation}
\end{equation}
No scene-wide bounding box is voxelized: the output contains the
warped-surface footprint and only the additional blocks reached by transported
free-space support, giving
$|\mathcal X^+|\le
n_b^3\!\left(27|\mathcal B_V^+|+|\mathcal B_{\rm fs}^+|\right)$.
For $x\in\mathcal X^+$, a closest-triangle
query~\cite{10.1145/2601097.2601199} on
$\mathcal M^+$ returns face $f$, closest point $q^+$, and barycentric
coordinates $\beta$. Shared face indices locate the corresponding source point
\begin{equation}
q^-=\sum_{k\in f}\beta_k\bar v_k^-.
\label{eq:source_correspondence}
\end{equation}
Let
\begin{equation}
R_f^\pm=
\begin{bmatrix}
t_{f,1}^\pm & t_{f,2}^\pm & n_f^\pm
\end{bmatrix}\in SO(3)
\label{eq:triangle_frames}
\end{equation}
be consistently oriented orthonormal frames of the corresponding source and
warped triangles. The first tangent follows the same indexed edge and
$t_{f,2}^\pm=n_f^\pm\times t_{f,1}^\pm$. A piecewise local-rigid pullback,
motivated by deformation transfer~\cite{sumner2004deformationtransfer}, maps
$x$ to the source volume:
\begin{equation}
\Psi_f(x)=q^-+R_f^-(R_f^+)^\top(x-q^+),
\qquad x^-=\Psi_f(x).
\label{eq:inverse_map}
\end{equation}
This preserves tangential and normal offsets relative to the corresponding
triangles.
With source TSDF samples $\widetilde\Phi^-$ and $\widetilde w^-$, let
\begin{equation}
r(x)=\|x-q^+\|,
\qquad
d^+(x)=\min\!\left(\tau,r(x)\right).
\label{eq:warped_distance}
\end{equation}
A complete trilinear source sample is valid only when every source voxel with
a nonzero interpolation coefficient carries positive weight. When such a
sample is available, WarpVox inherits its interpolated weight and its sign
evidence. When a complete distance sample is unavailable but the pullback
still reaches observed source support, only the source weight is inherited;
the corrected distance remains defined by the warped surface. Let
$w_{\rm src}^+(x)$ denote the resulting inherited source weight.

Let $s_f^+(x)$ denote the oriented warped-surface sign. Whenever possible,
the corresponding source TSDF orients the warped normal using its variation
along the source-triangle normal. If this orientation is unavailable, a
sign-consistent source sample is used in the surface-local band; otherwise
the oriented warped-surface sign is retained. We denote the resulting sign by
$s_{\rm tr}(x)$.

The initial source-supported candidates are
\begin{equation}
\begin{aligned}
\widehat{\Omega}_{\rm tr}
&=\left\{x\in\mathcal X^+:
r(x)\le\tau,\ w_{\rm src}^+(x)>0\right\},\\
\Omega_{\rm n}
&=\left\{x\in\widehat{\Omega}_{\rm tr}:r(x)<1.5s\right\},\\
\widehat{\Omega}_{\rm e}
&=\widehat{\Omega}_{\rm tr}\setminus\Omega_{\rm n}.
\end{aligned}
\label{eq:transferred_support}
\end{equation}

Let $\partial_{\rm open}\mathcal M^+$ denote the union of open boundary edges
and their incident vertices. Where source support is unavailable, unsupported
voxels in the same finite surface band receive a unit fallback prior:
\begin{equation}
\Omega_{\rm fb}=
\left\{x\in\mathcal X^+\setminus\widehat{\Omega}_{\rm tr}:
r(x)<1.5s,\ q^+\notin\partial_{\rm open}\mathcal M^+\right\}.
\label{eq:fallback_support}
\end{equation}
Its weight is
\begin{equation}
w_{\rm fb}(x)=
\operatorname{clamp}_{[0,1]}\!\left(
1.5-\frac{r(x)}{s}\right).
\label{eq:fallback_weight}
\end{equation}
This fallback is a finite prior required to serialize and continue updating
the corrected zero set; it is not interpreted as recovered observation count.

From $\widehat{\Omega}_{\rm e}$, WarpVox retains only candidates that are
sign-consistent with neighbouring active support and connected to the
surface-local support $\Omega_{\rm n}\cup\Omega_{\rm fb}$. Let the retained
set be $\Omega_{\rm e}$ and define
$\Omega_{\rm tr}=\Omega_{\rm n}\cup\Omega_{\rm e}$. For
$x\in\Omega_{\rm tr}$,
\begin{equation}
\Phi_{\rm tr}^+(x)=s_{\rm tr}(x)d^+(x),
\qquad
w_{\rm tr}^+(x)=w_{\rm src}^+(x).
\label{eq:transported_tsdf}
\end{equation}
To preserve observed free space beyond the surface-local bands, confidently
positive source TSDF voxels are grouped locally and transported forward using
the paired source and warped surface frames. For a local group $c$, let
$q_c^-$ and $q_c^+$ be corresponding points on the paired triangles and let
$R_c^-$ and $R_c^+$ be their local frames. The forward map is
\begin{equation}
\Gamma_c(y)=
q_c^+ + R_c^+(R_c^-)^\top(y-q_c^-).
\label{eq:forward_free_space}
\end{equation}
Mapped samples may allocate additional output blocks. They are written only
into previously unsupported voxels and are rejected when they would introduce
a new zero crossing adjacent to retained occupied support. The accepted
samples define $\Omega_{\rm fs}$ and retain positive, source-derived TSDF
values and weights $(\Phi_{\rm fs}^+,w_{\rm fs}^+)$.

Define
$\Omega^+=\Omega_{\rm tr}\cup\Omega_{\rm fb}\cup\Omega_{\rm fs}$.
The branches are disjoint because fallback and free-space transport write only
to unsupported voxels. The reconstructed TSDF is
\begin{equation}
\Phi^+(x)=
\begin{cases}
\Phi_{\rm tr}^+(x), & x\in\Omega_{\rm tr},\\
s_f^+(x)d^+(x), & x\in\Omega_{\rm fb},\\
\Phi_{\rm fs}^+(x), & x\in\Omega_{\rm fs},
\end{cases}
\label{eq:hybrid_tsdf}
\end{equation}
with weights
\begin{equation}
w^+(x)=
\begin{cases}
w_{\rm tr}^+(x), & x\in\Omega_{\rm tr},\\
w_{\rm fb}(x), & x\in\Omega_{\rm fb},\\
w_{\rm fs}^+(x), & x\in\Omega_{\rm fs},\\
0, & \text{otherwise}.
\end{cases}
\label{eq:hybrid_weight}
\end{equation}
Because each branch assigns positive weight on its support, by construction
$\Omega^+=\{x\in\mathcal X^+:w^+(x)>0\}$.

When no source TSDF is available, the final Stage-3 surface is converted into
a closest-triangle signed Euclidean distance field within $1.5$ voxel widths,
using the same finite fallback prior. The resulting $(\Phi^+,w^+)$ is
serialized in the native Voxblox TSDF format; its MC mesh is used for surface
evaluation and visualization. The returned volume replaces the active Voxblox
TSDF, so subsequent fusion and later correction events proceed from the
corrected state.

\subsection{Complexity and Scope} Let $A=|\mathcal A|$, $K_{\rm KF}=|\mathcal C_{\rm KF}|$, $P=|P^-|$, and $X^+=|\mathcal X^+|$ be the numbers of active source cells, stored keyframe range images, proxy vertices, and allocated output voxels. At fixed resolution, MC yields a constant number of surface elements per active cell, so $|V^-|=O(A)$; storage is $O(\sum_{i\in{\rm KF}}|D_i|)$ for the measurements and $O(A)$ for the per-vertex association state. With $\kappa\le K_{\rm KF}$ the maximum number of keyframes returned for one source vertex, association costs $O(K_{\rm KF}\log K_{\rm KF}+A(\log K_{\rm KF}+\kappa))$ for RGB-D and $O(K_{\rm KF}+A(1+\kappa))$ for LiDAR under the standard output-sensitive index model, giving \begin{equation} C_{\rm corr}=C_{\rm assoc}+C_{\rm ARAP}(P,F_P)+C_{\rm vol}(X^+,\bar F), \label{eq:complexity} \end{equation} where $C_{\rm vol}$ is the Stage-4 volumetric reconstruction cost. The block-local allocation in \eqref{eq:allocation} makes $X^+$ follow the warped-surface footprint and transported source support rather than scene extent, and the proxy targets approximately 60K vertices. The retained association state per surface element is bounded by a constant, and each output voxel is written once rather than once per contributing measurement. Full reintegration instead re-fuses every historical measurement at $O(N|D|)$. The two costs are indexed by different quantities: the keyframes and the active map for WarpVox, the measurement history for reintegration.
